% !TEX root = ../main.tex
% CHU SO HUU: Mỹ Uyên và Mai Quyên | KHONG THEO PREAMBLE, KHONG BEGIN DOCUMENT TRONG THAN PHAN NAY.
\documentclass[../main.tex]{subfiles}
\begin{document}
\section{Phân tích tương quan Pearson}
\label{sec:t31:phuong-phap}

NSL-KDD được xây dựng nhằm cải thiện một số hạn chế của KDD Cup 99
\cite{task1_tavallaee2009}. Các đặc trưng thời lượng, lưu lượng và trạng thái
kết nối có thể cùng phản ánh một hành vi mạng, tạo thông tin tuyến tính
trùng lặp. Nhóm dùng tương quan Pearson để xem xét khả năng dư thừa,
không tự động loại biến dựa trên tương quan cao.

\subsection{Hệ số Pearson và quy trình lựa chọn biến}
\label{sec:t31:pearson-r}
Hệ số tương quan Pearson mẫu giữa hai biến định lượng $X,Y$ được tính bởi
\begin{equation}
 r_{XY}=\frac{\sum_{i=1}^{n}(x_i-\bar{x})(y_i-\bar{y})}
 {\sqrt{\sum_{i=1}^{n}(x_i-\bar{x})^2}\sqrt{\sum_{i=1}^{n}(y_i-\bar{y})^2}}.
 \label{eq:t31:pearson}
\end{equation}
Trong đó $n$ là số kết nối, $x_i,y_i$ là giá trị quan sát thứ $i$,
còn $\bar{x},\bar{y}$ là các trung bình mẫu. Với $r\in[-1,1]$,
dấu dương/âm thể hiện chiều tương quan; $|r|$ càng gần 1 thì tương quan
tuyến tính càng mạnh. $r$ gần 0 không loại trừ quan hệ phi tuyến;
tương quan không chứng minh quan hệ nhân quả.

Phân tích dùng $125\,973$ kết nối của tập huấn luyện \texttt{KDDTrain+},
nhất quán với \texttt{train\_clean.csv}. Loại ba biến danh nghĩa
\texttt{protocol\_type}, \texttt{service}, \texttt{flag}, các cột
nhãn/metadata và biến hằng \texttt{num\_outbound\_cmds}. Còn 37 đặc trưng
định lượng, xét các cặp không trùng lặp; ngưỡng $|r|>0{,}95$ dùng để sàng lọc
cặp tương quan rất cao, không phải quy tắc tự động loại biến.

Ba biến danh nghĩa nêu trên không được đưa trực tiếp vào tính Pearson vì mã
số nhóm không biểu thị khoảng cách định lượng. Bốn cột nhãn/metadata được
loại khỏi tập đặc trưng gồm \nolinkurl{attack_name}, \texttt{difficulty},
\nolinkurl{attack_category} và \nolinkurl{label_binary}. Biến
\texttt{num\_outbound\_cmds} có phương sai bằng không nên không thể
tính hệ số Pearson với các biến còn lại.

\subsection{Kiểm định ý nghĩa thống kê của tương quan}
\label{sec:t31:kiem-dinh}
Nhóm kiểm định hai phía trên chín cặp tiêu biểu cho tương quan rất cao,
trung bình, âm, yếu và gần bằng không. Với mỗi cặp $X,Y$, đặt
\begin{equation}
 H_0:\rho_{XY}=0,\qquad H_1:\rho_{XY}\neq 0,
 \label{eq:t31:hypotheses}
\end{equation}
trong đó $\rho_{XY}$ là hệ số tương quan Pearson của tổng thể, phân biệt
với hệ số mẫu $r_{XY}$. Chọn mức ý nghĩa $\alpha=0{,}05$. Thống kê kiểm định là
\begin{equation}
 t= r\sqrt{\frac{n-2}{1-r^2}}, \qquad df=n-2=125\,971.
 \label{eq:t31:t-stat}
\end{equation}
Giả sử các cặp quan sát độc lập và phân phối chuẩn hai chiều phù hợp,
dưới $H_0$ thì thống kê có phân phối Student với $df=n-2$ bậc tự do.
Giá trị $p$ hai phía được tính là $2P(T_{df}\geq|t_{\mathrm{obs}}|)$.
Quy tắc ra quyết định: bác bỏ $H_0$ khi $p<0{,}05$; nếu không,
\emph{chưa đủ bằng chứng} để bác bỏ $H_0$.

Vì lựa chọn cặp sau khảo sát ma trận, đây là phân tích \emph{khám phá}.
Nhóm tính thêm $p$ điều chỉnh Holm cho chín cặp, không phải cả 666 cặp.
Kết nối mạng có thể phụ thuộc thời gian/nguồn phát; biến đếm thường lệch và
có ngoại lệ, khiến giả định kiểm định chưa chắc thoả mãn. Vì vậy, cần đọc
$p$ cùng $r$, độ nhạy với ngoại lệ và ý nghĩa chuyên ngành; không tự động
suy rộng sang mọi lưu lượng mạng thực tế.

\end{document}
